%%%%%%%%%%  scan idx 65  |  folio 52  %%%%%%%%%%
par une
formule faisant intervenir une somme alternée de $\Tor$, inspirée par les développements du début du \S 2, et que le lecteur explicitera. On a alors commutativité dans le diagramme suivant~:
\[
\begin{tikzcd}[column sep=large,row sep=large]
\underline K(\xi(Y)) \arrow[r,"f^!"] \arrow[d] & \underline K(\xi(X)) \arrow[d]\\
\underline K(Y) \arrow[r,"f^!"] & \underline K(X){,}
\end{tikzcd}
\]
ce qui montre en particulier, lorsque $X$ et $Y$ sont toutes deux non singulières et quasi-projectives, que les deux définitions de $f^!$ coïncident moyennant l'identification de $\underline K(X)$ à $\underline K(\xi(X))$ et de $\underline K(Y)$ à $\underline K(\xi(Y))$.

Un autre cas où l'on peut définir un homomorphisme (2.21), sans que $Y$ soit nécessairement non singulière, est celui où le foncteur $F\to f^!(F)$ de $\cF(Y)$ dans $\cF(X)$ est exact, i.e.\ lorsque les anneaux locaux de $X$ sont des modules \emph{plats} sur les anneaux locaux de $Y$. Il suffit alors de poser $f^!(\gamma_Y(F))=\gamma_X(f^!(F))$.
Ce cas se rencontre en particulier lorsque $X$ est un espace fibré localement trivial sur $Y$ (pour le voir, on est ramené au cas d'un produit $X=T\times Y$). Notons qu'on vérifie aussi que dans ce cas, on a $f^!(\cO_Z)=\cO_{f^{-1}(Z)}$ si $Z$ est une sous-variété irréductible de $Y$, d'où~:
\begin{equation}\tag{2.22}
 f^!(\gamma_Y(Z))=\gamma_X(f^{-1}(Z))
\end{equation}
\emph{si $X$ est «~plat~» sur $Y$}. Sans cette dernière condition, la formule (2.22) devient inexacte en général, mais on a l'analogue suivant de la proposition 2.7~:

\textbf{Proposition 2.8.}
\emph{Soit $f:X\longrightarrow Y$ un morphisme de variétés non singulières, et soit $Z$ un cycle sur $Y$. Si $Z$ est partout transversal à $f$, la formule}
